\documentclass[11pt]{article}
\usepackage[margin=1in]{geometry}
\usepackage{amsmath,amssymb,amsthm,mathtools}
\usepackage[T1]{fontenc}
\usepackage{hyperref}
\newtheorem{theorem}{Theorem}
\newtheorem{lemma}[theorem]{Lemma}
\newtheorem{proposition}[theorem]{Proposition}
\newtheorem{corollary}[theorem]{Corollary}
\theoremstyle{remark}
\newtheorem{remark}[theorem]{Remark}
\newcommand{\R}{\mathbb R}
\newcommand{\Z}{\mathbb Z}
\newcommand{\T}{\mathcal T}
\newcommand{\Hs}{H^s}
\newcommand{\Wh}{\operatorname{Wh}^{\mathrm{Diff}}}
\newcommand{\Crit}{\operatorname{Crit}}
\newcommand{\im}{\operatorname{im}}
\newcommand{\norm}[1]{\left\lVert #1\right\rVert}
\newcommand{\ip}[2]{\left\langle #1,#2\right\rangle}
\title{Nonzero stable tubes and their exact embedded realization\\
\large A bounded verification of the construction in manuscript 340}
\author{Campaign 9000340}
\date{5 November 2026}
\begin{document}
\maketitle

\begin{abstract}
We verify the closed exact embedded realization in the manuscript's
tube and embedding sections. We specify the component,
finite-index shift, spherical classifier, and loop convention in the
imported tube theorem. The actual connecting homomorphism remains
unknown: an exhaustive finite-group certificate covers every possible
homomorphism with the stated constraints, not a purported numerical
representative of the topological class. The construction is proved for
all points and parameters, including singular vertical Hessians and
collision segments through the origin. No Hamiltonian-transport theorem
or novelty assertion is made.
\end{abstract}

\section{Statement and inputs}

Put \(q_{k,l}(x,y)=-|x|^2+|y|^2\). Let \(T_{k,l}\) be the \(C^1\)
path component of \(q_{k,l}|_{S^{k+l-1}}\) in the angular functions having
zero as a regular value, and set
\[
 \T=\operatorname*{hocolim}_{m}
       (T_{m,m}\xrightarrow{\oplus q_{1,1}}T_{m+1,m+1}).
\]
For \(h\) in this component write
\[
 g(ru)=r^2h(u)\quad(r>0),\qquad g(0)=0,\qquad
 D(h)=\{h\leq0\}\subset S^{k+l-1}.
\]
The map \(\gamma:\T\to BG\) classifies the stable spherical fibration
of these \emph{negative domains}. It is not the gradient or
twisted-derivative normal invariant.

\begin{theorem}\label{thm:main}
For some even \(N\geq40\), there is a smooth function
\[
 F:(S^9\times S^{N-1})\times\R^N\longrightarrow\R
\]
with the following properties. Zero is a regular value of its full
vertical-gradient map. Its critical locus is diffeomorphic to
\(S^9\times S^{N-1}\), and
\[
 \Gamma_F(b,v,w)=(b,v,d_{b,v}F)
\]
embeds that locus as a closed exact Lagrangian in
\(T^*(S^9\times S^{N-1})\). Outside one uniform vertical ball,
\(F(b,v,w)=g_b(w)\), independent of \(v\), and on every
\(S^9\times\{v_0\}\) its stable tube class is the same nonzero
\(a\in\ker\{\pi_9\T\to\pi_9BG\}\).
\end{theorem}

The topological inputs are stated explicitly here. They are imported
theorems, not conclusions of the computational certificate.
\begin{enumerate}
\item The stable smooth tube fibration and rigid map are
\[
 \Hs(*)\longrightarrow\T\xrightarrow{\gamma}BG,\qquad
 BO\xrightarrow{\iota}\T,\qquad \gamma\iota\simeq BJ.
 \tag{T}
\]
Waldhausen's original Propositions 3.1--3.2 and Corollary 3.4
\cite{Waldhausen} give the partition model. The spherical-domain
comparison, Theorem 3.30, and Remark 3.34 in the exact version
\cite{AGIS} supply its smooth tube formulation.
Definition 4.3 and Propositions 4.4, 4.6--4.8 of \cite{AAGCK}
identify the negative-domain classifier. Section~\ref{sec:model}
checks that these statements concern the model used here.
\item The stable smooth \(h\)-cobordism theorem gives
\[
 \Hs(X)\simeq\Omega\Wh(X),\qquad
 \mathbf{Wh}(X)=\operatorname{cofib}
        \{\Sigma^\infty X_+\longrightarrow\mathbf A(X)\}.
 \tag{W}
\]
Here \(\Wh(X)=\Omega^\infty\mathbf{Wh}(X)\).
Theorems 0.1 and 0.3 of the retained author-hosted WJR text
\cite{WJR} have precisely this one-loop convention. That inspected
file has 2012 prepublication metadata; we do not claim a separate
inspection of the final 2013 edition.
\item Adams's Theorem 9.5 \cite{Adams} states that
\(e_{\R}J:\pi_r(SO)\to\Z/2\) is an isomorphism for
\(r>1\), \(r\equiv0,1\pmod8\). In particular both \(J_8\) and \(J_9\)
are injective with order-two images. No additional use of Bott
periodicity is necessary.
\item The published Table 1, degree 9, of \cite{IWX}, with its
column explanation and Corollary 1.3, gives
\(\pi_9^S=(\Z/2)^3\). The three direct-sum constituents are
\(\Z/2\), zero, and \((\Z/2)^2\); the smooth-structures column is
not another constituent. The table is numbered \emph{1}, not 1.1.
Rognes's Theorem 3.3 in \cite{Rognes} gives
\[
  \pi_{10}\mathbf{Wh}(*)_{(2)}
       \cong \Z/8\oplus(\Z/2)^2.
 \tag{G}
\]
Its qualification about torsion at irregular primes does not affect
this two-primary assertion.
\end{enumerate}
The relevant source passages were read by a separate source researcher;
the originals are retained in \texttt{sources/}, with Adams and WJR
retained in the earlier campaign directories. This is not a new
independent review of the present proof. The accepted full one-critical
obstruction is a separate, already reviewed campaign result. The
construction below neither reopens nor uses it. Its use to obstruct a
Hamiltonian isotopy still requires the separate transport argument.

\section{The actual tube model}\label{sec:model}

\subsection{Regularity, smooth families and domains}
For \(r>0\),
\begin{equation}\label{eq:angular}
 \nabla g(ru)=r(2h(u)u+\nabla_Sh(u)),\qquad
 |\nabla g(ru)|^2=r^2(4h(u)^2+|\nabla_Sh(u)|^2).
\end{equation}
The two summands in the first expression are orthogonal. Thus the
regular-zero condition is equivalent to \(\nabla g(w)\ne0\) for \(w\ne0\).
Also \(g=O(|w|^2)\), \(dg=O(|w|)\), so its extension at zero is \(C^1\).
This matches the meaning of ``quadratic'' in Definitions 3.10 and 3.36
of \cite{AAGCK}: positive homogeneity of degree two, not a requirement
that \(g\) be a polynomial or even under \(w\mapsto-w\).

A compact family of regular angular functions has a positive minimum
\[
 \mu=\min_{b,u}(4h_b(u)^2+|\nabla_Sh_b(u)|^2)^{1/2}.
\]
An angular \(C^1\) perturbation whose corresponding vector
\((2\Delta h,\nabla_S\Delta h)\) has norm less than \(\mu/2\)
preserves regularity throughout straight interpolation. A continuous
map from a compact smooth parameter manifold into \(C^1\) can be
approximated uniformly in this norm by finitely many smooth angular
functions with smooth parameter partition-of-unity coefficients.
If the family is already the standard form on an open parameter set,
multiply the approximating difference from the standard form by a
parameter cutoff supported away from a smaller such set.
The same construction for homotopies, relative to their ends, proves
that smooth compact families compute the homotopy groups of the
specified \(C^1\) component.

For a fixed smooth cooriented domain \(D\), defining angular functions
are negative in its interior, positive outside, and have positive
outward derivative on its boundary. These three conditions persist
under convex combinations. A boundary collar supplies one such
function. Nearby boundaries are normal graphs; their graph isotopies
extend in the collar, giving local choices of defining functions.
Over a compact parameter space, a partition of unity combines these
choices because, at each parameter, all summands define the
\emph{same} domain. This also extends prescribed endpoint functions:
choose their collar extensions first and interpolate within the same
convex fibers. Therefore passage between smooth defining functions
and smooth embedded, cooriented negative domains preserves compact
families and relative homotopies. No normal framing is included.

There is a useful exact check against Definition 4.1 of \cite{AAGCK}.
That definition requires \emph{existence} of a homogeneous ambient
isotopy over a simplex with \(f\circ\varphi_1=q\); an isotopy is not
extra simplex data. On a contractible parameter simplex, concatenate
a path from \(q\) to one vertex with the contraction paths inside the
family. For a smooth such path \(h_t\), the vector field
\[
 -\frac{\partial_t h_t}{|\nabla_Sh_t|^2}\nabla_Sh_t
\]
near its zero level, extended by a cutoff, integrates to sphere
diffeomorphisms \(\psi_t\) carrying the standard zero set and its
negative side to those of \(h_t\). The construction is uniform in
the simplex parameters. Having the same cooriented regular zero set,
\(q(u)\) and \(h_t(\psi_t(u))\) have a positive smooth ratio, including
across that set: in a transverse coordinate both are that coordinate
times a nonvanishing coefficient of the same sign. Consequently
\[
 \varphi_t(ru)=r
  \sqrt{\frac{q(u)}{h_t(\psi_t(u))}}\,\psi_t(u)
\]
is a homogeneous homeomorphism of the whole vector space and satisfies
\(g_t\circ\varphi_t=q\). It is smooth off zero and continuous at zero,
with continuous inverse. The source's normalization \(q/2\) is
handled by replacing \(g_t\) by \(g_t/2\).
Face compatibility is a condition on the functions, not on a chosen
isotopy. In particular this argument supplies the required isotopies
over simplices; it does \emph{not} choose one over all of \(S^9\),
which would erroneously trivialize the class being constructed.

\subsection{Stabilization, the standard-ball condition and the rigid map}
The direct sum of two homogeneous \(C^1\) functions has block
gradient, so it has no nonzero critical point. It may fail to be
smooth on a coordinate axis. Angular smoothing with the preceding
regularity margin removes this defect without changing its class.
Adding positive squares has negative domain, in the chart
\(u=w/|w|,\ z'=z/|w|\),
\[
             h(u)+|z'|^2\leq0.
\]
It is the rounded disk thickening of \(D(h)\). Adding negative
squares is the complementary construction and suspends the negative
spherical core. The two additions commute after the fixed coordinate
identifications. The diagonal is cofinal, so every finite combination
of the two operations is seen in the balanced telescope.

The source \cite{AAGCK} uses left direct sum, whereas one may write
the manuscript stabilization on the right. Reversing the order of
the coordinates within each negative and positive block gives an
exact isomorphism of these directed systems: the reversal of
\((\text{new},x_1,\ldots,x_k)\) is
\((x_k,\ldots,x_1,\text{new})\), in both blocks. These coordinate
maps fix the standard form and identify its negative eigenspace
bundle by a fixed linear isometry. Thus this convention change
does not replace \(\gamma\) or the rigid stable \(J\)-map.

We spell out why the extra standard-ball condition in the partition
model is harmless \emph{stably}, including its relative meaning.
Use the collared standard-ball model corresponding to partitions
standard outside a compact set. A compact family in that model has
a common standard collar; no assertion that agreement on a closed
ball alone supplies a larger collar is needed.
Take a smooth family over a compact \(d\)-dimensional parameter
manifold \(P\), possibly with boundary. The boundary-domain bundle
has fiber \(S^{k-1}\times S^{l-1}\), hence connectivity
\(\min(k,l)-2\). After increasing both indices, obstruction theory
supplies a section, also extending a prescribed section over
\(\partial P\). Its image in the ambient sphere can be homotoped
to a fixed point, relative to that boundary, once \(k+l-1>d\).
Parametric isotopy extension moves the boundary points accordingly.
The outward unit normal at the fixed point is then a map to
\(S^{k+l-2}\), which likewise becomes relatively nullhomotopic.
An ambient isotopy fixing the point makes all these normals fixed.
For example \(k,l\geq d+3\) safely suffices for all these steps.

First make parameter collars stationary by pulling back along a
collar retraction, relative to the parameter boundary. In the
resulting uniform chart the domain boundaries are graphs over a
fixed hyperplane. Compactness supplies a common graph neighborhood.
Flattening these graphs makes the domains the same half-space on a
small ball. A fixed hyperplane-preserving radial isotopy enlarges
this small standard patch to the standard ball of the partition
model. For the relative assertion, on the prescribed parameter
boundary the domains already agree on that ball and a collar;
compactness gives a common slightly larger collar. The enlargement
is supported there and preserves its standard half-space, so it
does not change those boundary domains. Earlier straightening steps
can also be fixed there. These arguments apply to \(P=S^j\) and to
relative nullhomotopies on \(D^{j+1}\), and prove precisely the
stable, not a finite-stage, comparison used in (T). They are the
standard-ball comparison of \cite[Section 3.6]{AGIS}.

One must compare finite indices by geometry, not by repeated symbols.
Our negative domain has core \(S^{k-1}\) and ambient dimension
\(k+l-1\). Waldhausen's \(T^{a,b}\), in ambient
\(\R^{a+b+1}\), has core \(S^a\). Thus the correspondence is
\begin{equation}\label{eq:index-shift}
                 a=k-1,\qquad b=l-1,
\end{equation}
not \(a=k,b=l\). For a rigid tube, choosing the boundary point and
its normal splits off one negative and one positive line. Indeed if
\(a_-\) and \(a_+\) are unit vectors in those eigenspaces, the point
and outward normal are
\((a_-+a_+)/\sqrt2\) and \((-a_-+a_+)/\sqrt2\).
The remaining subspaces have dimensions \(k-1,l-1\).
Both indices increase by one under \(q_{1,1}\); the finite shift
therefore respects the directed systems and does not change stable
spherical data. Ambient isotopies preserve the inherited ambient
stable structure; an arbitrary normal framing has not been supplied.

There are also component labels in the sources. The stabilized whole
disjoint union in \cite{AAGCK} has components indexed by
\((k-i,l-i)\in\Z^2\), at telescope stage \(i\). Its printed coordinates
\((k-l,k+l-2i)\) have the parity relation that their two entries are
congruent modulo two; they describe the same rank-two lattice.
Our diagonal selects a single component, not all of
\(\Z^2\times BG\). A different initial stage just translates its
label. Increasing both indices, or passing to the cofinal
\(q_{2,2}\) subsystem, does not alter this component's stable
homotopy type. The latter observation also avoids confusing the
sign change of higher torsion under one index suspension with a
failure of the underlying tube comparison. No higher-torsion
formula is used in this proof.

In Definition 4.3 of \cite{AAGCK}, a negative frame is a homogeneous
map \(\R^k\setminus0\to\{g<0\}\) that is a homotopy equivalence.
Its automorphism monoid is the unbased homotopy self-equivalence
monoid of \(S^{k-1}\). The open cone retracts radially to its angular
part, and the regular boundary collar identifies the homotopy types
of its open and closed negative angular domains. Thus the associated
classifier really is \(\gamma\). Proposition 4.4 identifies the
finite-stage bar construction \(B(E_{k,l};G_k^u)\), not the finite
stage with \(BG_k^u\); Proposition 4.6 makes the frame space
contractible only after both indices increase. Proposition 4.7 then
gives \(BG\) on the component in question.
The evaluation fibration has a sphere of increasing connectivity as
base, so based and unbased \(G\) agree in the stable positive
homotopy groups used here. Both degree signs are retained.

For a quadratic form the negative domain retracts onto the sphere
of its negative eigenspace. The balanced Grassmannian telescope
is the index-zero component \(BO\), and this construction sends
it to the stable sphere bundle of that eigenspace, namely \(BJ\).
This is also the negative-linear-frame comparison in
\cite[Proposition 4.8]{AAGCK}. It proves the compatibility of the
rigid map in (T) with the actual domain classifier. Together with
the imported stable smooth tube theorem, the model check is complete.

\section{The kernel, without inventing its boundary map}
The relevant exact sequence is
\begin{equation}\label{eq:les}
 \pi_{10}\T\longrightarrow\pi_{10}BG
 \xrightarrow{\partial}\pi_9\Hs(*)
 \xrightarrow{j_*}\pi_9\T\xrightarrow{\gamma_*}\pi_9BG.
\end{equation}
The shifts are
\[
 \pi_{10}BG=\pi_9G=\pi_9^S,\quad
 \pi_9\Hs(*)=\pi_{10}\mathbf{Wh}(*),\quad
 \pi_{10}BO=\pi_9O,\quad \pi_9BO=\pi_8O.
\]
Thus the lifting rigid map puts the order-two image of \(J_9\)
in \(\ker\partial\); injectivity of \(BJ\) in degree nine instead
uses \(J_8\). These are different degrees.

Localize the exact sequence at two. Localization is exact for these
abelian groups. Write \(A=(\Z/2)^3\), \(H=(\Z/2)^2\oplus\Z/8\).
The source is already two-primary. Since its order-two rigid
subgroup lies in the kernel, \(|\im\partial|\leq8/2=4\), whereas
\(|H|=32\). Consequently
\begin{equation}\label{eq:kernel-bound}
 H/\im\partial\ \lhook\joinrel\longrightarrow\
       \ker(\gamma_*)_{(2)},\qquad |H/\im\partial|\geq8.
\end{equation}
In particular the unlocalized kernel is nonzero. No surjectivity
or formula for \(\partial\) has been asserted.

\begin{proposition}\label{prop:universal}
For every homomorphism \(A\to H\) annihilating a specified nonzero
vector, the possible cokernels, and numbers of such homomorphisms
with each cokernel, are
\[
\begin{array}{c|r}
\text{cokernel}&\text{number}\\ \hline
 (\Z/2)^2\oplus\Z/8&1\\
 \Z/2\oplus\Z/8&18\\
 (\Z/2)^2\oplus\Z/4&3\\
 \Z/8&24\\
 \Z/2\oplus\Z/4&18
\end{array}
\]
Every such cokernel has a canonical quotient
\(H/H[2]\cong\Z/4\). This table does not select which case is the
topological connecting map.
\end{proposition}
\begin{proof}
Choose a source basis with the killed vector first. The images of
the other two basis vectors may be any ordered pair in
\[
 H[2]=(\Z/2)^2\oplus\{0,4\},
\]
which has eight elements. Thus there are exactly \(8^2=64\) maps.
Their image dimensions over \(\mathbb F_2\) are zero, one or two.
There is one rank-zero map. Each of the seven lines is the image
of three maps; each of the seven planes is the image of six maps.

Let \(z=(0,0,4)\). A rank-one image containing \(z\) gives
\((\Z/2)^2\oplus\Z/4\); the other six lines give
\(\Z/2\oplus\Z/8\). Among the seven planes three contain \(z\)
and yield \(\Z/2\oplus\Z/4\); the other four yield \(\Z/8\).
For example these types are also determined directly by the orders
of elements in the quotient. This gives the table.
Every image is contained in \(H[2]\), so passage to the larger
subgroup gives the asserted quotient \(\Z/4\).
\end{proof}

The standalone program independently builds every quotient as actual
cosets and checks all addition, kernel, element-order and type counts.
The retained Jacobian call computes only the canonical
\(H/H[2]\), with generators \((1,0,0),(0,1,0),(0,0,4)\).
Those generators are \emph{not} an asserted image of \(\partial\).
This also shows that nonzero kernel existence itself does not need
the sharper \(J_9\) order bound: an exponent-two image cannot exhaust
the order-eight summand. We nevertheless checked \(J_9\), as used
by the manuscript.

\section{Finite-stage localization}
\begin{lemma}\label{lem:localized}
For arbitrarily large \(m\geq20\), a nonzero class in
\(\ker\gamma_*\subset\pi_9\T\) has a smooth representative
\(b\mapsto h_b\in T_{m,m}\), equal to the standard form on
neighborhoods of both critical points of a standard height
\(\theta:S^9\to[-1,1]\). Its extension satisfies
\begin{equation}\label{eq:bounds}
 c|w|\leq|\nabla g_b(w)|\leq C|w|,\qquad
 \norm{d_wd_bg_b(w)}\leq C_1|w|
\end{equation}
for uniform \(c>0\), \(C<\infty\), \(C_1<\infty\).
The covector-valued function \(w\mapsto d_bg_b(w)\) is \(C^1\)
at zero, with derivative zero.
\end{lemma}
\begin{proof}
The telescope has a continuous height coordinate. A compact sphere
map has bounded height, and a finite initial telescope retracts onto
its last stage. Hence the class has a finite-stage representative;
further stabilization makes \(m\) arbitrarily large. Smooth the
angular functions and parameters as in Section~\ref{sec:model}.

Choose a small oriented disk disjoint from both height-critical
points. A based degree-one collapse onto \(S^9\), constant off that
disk, followed by relative smoothing gives a smooth degree-one
self-map constant on neighborhoods of the two points. A based
degree-one self-map of the sphere is based homotopic to the identity.
Precomposition therefore preserves both the chosen class and its
spherical triviality. Relative angular approximation retains the
constant standard functions on smaller neighborhoods.

Compactness of \(S^9\times S^{2m-1}\) in \eqref{eq:angular} gives
the first two bounds. In local base coordinates write
\(d_bg_b(ru)=r^2a_b(u)\), where \(a_b=d_bh_b\) is a smooth
covector-valued angular function. Its \(w\)-derivative is
\[
 r\bigl(2a_b(u)u^\flat+d_Sa_b(u)\bigr),
\]
interpreted as a linear map on the radial/tangential decomposition.
Compactness bounds its norm by \(C_1r\).
The value is \(O(r^2)\), so its derivative at zero exists and is zero;
the \(O(r)\) derivative estimate proves continuity of that derivative.
All estimates are uniform in base charts and hence in a fixed base
metric. A nonzero based sphere class is not freely nullhomotopic,
because the fundamental-group action fixes zero.
\end{proof}

\section{The global embedding}
We give all uniform choices and both momentum tests.
Let \(B=S^9\), \(N=2m\), and let \(g_b\) be the localized family.
Choose an open union \(U_0\) of pole neighborhoods on which
\(g_b=q_{m,m}\), small enough that
\(\delta_0=\min_{B\setminus U_0}|d\theta|>0\).
Fix a smooth cutoff \(\chi(r)\) equal to zero near zero and to one
for \(r\geq1\), and put \(G_b(w)=\chi(|w|)g_b(w)\), extended by
zero near the origin. This is a smooth joint family. Define
\[
 M=\sup_{b,|w|\leq1}|\nabla G_b(w)|<\infty.
\]
No estimate for \(g_b\) at its nonsmooth second derivative at zero
is being used for \(G_b\).

Let \(0\leq\beta\leq1\) be smooth, equal to one on \([0,1]\)
and zero on \([2,\infty)\), and let
\[
 D=\sup_{s\geq0}\bigl(|\beta(s)|+s|\beta'(s)|\bigr).
\]
In particular \(D\geq1\). There is no need to specify a numerical
cutoff for the unknown topological representative. The following
choices are explicit in the family-dependent bounds:
\begin{equation}\label{eq:choices}
 \begin{aligned}
 K&=1+\frac{C_1}{c\delta_0},&
 \Lambda&=(1+C+M)e^K,\\
 R(b)&=\Lambda e^{K\theta(b)},&
 R_{\min}&=1+C+M,\\
 R_{\max}&=(1+C+M)e^{2K},&
 T&=1+(1+D)R_{\max}/c.
 \end{aligned}
\end{equation}
Since the height reaches \(\pm1\), these minimum and maximum
formulas are exact. They give
\[
 K\delta_0-C_1/c=\delta_0>0,\quad
 R_{\min}>C,M,\quad
 T>1,R_{\max}/c,DR_{\max}/c.
\]
The order is \(K,\Lambda,T\); none of these bounds depends on \(T\)
through a previously chosen cutoff.
Set
\begin{equation}\label{eq:F}
 F(b,v,w)=G_b(w)-R(b)\beta(|w|/T)\ip{v}{w},
       \qquad v\in S^{N-1}.
\end{equation}
It is smooth also at \(w=0\), since both radial cutoffs are constant
there.

\subsection{Every critical point, including the cutoff regions}
When \(|w|\leq1\), the cutoff \(\beta\) is constant and
\(|d_wF|\geq R_{\min}-M>0\).
When \(r=|w|\geq T\), \(G_b=g_b\), and
\[
 d_w\{\beta(r/T)\ip{v}{w}\}
  =\beta(r/T)v^\flat+
       \frac{\beta'(r/T)}{T}\ip{v}{w}\frac{w^\flat}{r}
\]
has norm at most \(|\beta|+(r/T)|\beta'|\leq D\).
It follows that \(|d_wF|\geq cr-DR(b)>0\).
This covers both annulus endpoints, the entire transition \(T\leq
r\leq2T\), and all larger \(r\).

In the remaining region \(1<r<T\), criticality is exactly
\begin{equation}\label{eq:critical}
                     \nabla g_b(w)=R(b)v.
\end{equation}
Conversely every solution of this equation satisfies
\[
           1<R(b)/C\leq |w|\leq R(b)/c<T,
\]
so it really lies in the unchanged region and is a critical point
of \eqref{eq:F}. With \(A_b(u)=\nabla g_b(u)\), the whole critical
locus is consequently parametrized by the diffeomorphism
\begin{equation}\label{eq:param}
 (b,u)\longmapsto
 \left(b,\frac{A_b(u)}{|A_b(u)|},
                 \frac{R(b)u}{|A_b(u)|}\right).
\end{equation}
Its inverse is \((b,v,w)\mapsto(b,w/|w|)\), smooth on this locus.
This identifies its \emph{abstract manifold} with the base
\(B\times S^{N-1}\); it does not assert that its projection to
\((b,v)\) is a diffeomorphism.

\subsection{Full Morse-family regularity, without Hessian invertibility}
At a critical point write \(H=D_w^2g_b(w)\). Homogeneity gives
\(Hw=\nabla g_b(w)=Rv\). For any target vector \(z\in\R^N\), take
\begin{equation}\label{eq:right-inverse}
 \delta b=0,\qquad
 \delta v=-\frac{z-\ip{v}{z}v}{R},\qquad
 \delta w=\frac{\ip{v}{z}}{R}\,w.
\end{equation}
Here \(\delta v\perp v\), so this is a legitimate tangent variation
in the sphere factor. The differential of \(d_wF\) on this variation is
\[
             H\delta w-R\delta v=z.
\]
Thus zero is a regular value of the \emph{full} gradient map,
even when \(H\) is singular.

For completeness the generating-map criterion is elementary.
In local coordinates the defining differential is the surjective
matrix \(A=(F_{wq},F_{ww})\), where \(q=(b,v)\).
If a tangent vector is killed by \(d\Gamma_F\), then
\(\delta q=0\), \(F_{ww}\delta w=0\), and
\(F_{qw}\delta w=0\). Symmetry of \(F_{ww}\) says
\(A^{\mathsf T}\delta w=0\), so \(\delta w=0\).
Thus \(\Gamma_F\) is an immersion. With
\(\lambda=p\,dq\) and \(\omega=d\lambda\),
\[
        \Gamma_F^*\lambda=d(F|_{\Crit(F)}).
\]
The critical locus has dimension \(\dim q\), so this exact
isotropic immersion is Lagrangian. Its primitive has sign
\begin{equation}\label{eq:primitive}
 F|_{\Crit(F)}=g_b(w)-R\ip{v}{w}=-g_b(w),
\end{equation}
by Euler's identity \(\ip{\nabla g_b(w)}{w}=2g_b(w)\).

\subsection{Both momentum components, for every collision pair}
If two critical points have equal images, their base coordinates
are first equal, say \((b,v)\). Their sphere covectors are the
restrictions of \(-Rw^\flat\) and \(-Rw'^\flat\) to \(v^\perp\).
Equality means
\begin{equation}\label{eq:parallel}
                       w-w'=tv
\end{equation}
for a real number \(t\). No other ambiguity remains after this test.
Both endpoints lie in the closed ball of radius \(R/c\), hence so
does their entire segment. The \(C^1\) extension at zero in
Lemma~\ref{lem:localized} makes the fundamental theorem of calculus
valid on that segment even if it passes through zero. It gives
\begin{equation}\label{eq:mixed}
 \begin{split}
 |d_bg_b(w)-d_bg_b(w')|
 &\leq\int_0^1 C_1|w'+s(w-w')|\,|w-w'|\,ds\\
 &\leq (C_1/c)R|t|.
 \end{split}
\end{equation}
This integral uses the reference \(g_b\), not the smoothed \(G_b\):
only its endpoint values are being identified with momenta of \(F\).
Passing through the origin-smoothing region therefore introduces
no omitted cutoff derivative into this identity.
At these points all cutoffs are constant. The difference of their
base covectors is therefore exactly
\[
 d_bg_b(w)-d_bg_b(w')-t\,dR
      =d_bg_b(w)-d_bg_b(w')-KRt\,d\theta.
\]
For \(b\notin U_0\), reverse triangle inequality and
\eqref{eq:choices} bound its norm below by
\[
             R|t|(K\delta_0-C_1/c)=R|t|\delta_0.
\]
Equality of these covectors forces \(t=0\), hence \(w=w'\).
For \(b\in U_0\), \(g_b=q_{m,m}\), and
\(\nabla q_{m,m}(w)=Rv\) has exactly one solution because its
constant Hessian is invertible. This includes both critical points
of the height and every point where the lower-bound argument was
not used.

The parametrization \eqref{eq:param} has compact domain without
boundary. Its image under the injective immersion \(\Gamma_F\)
is therefore an embedded, compact, boundaryless exact Lagrangian.
For \(|w|\geq2T\), \eqref{eq:F} is exactly \(g_b(w)\), independent
of \(v\). The class on each fixed sphere slice is unchanged, with
its given stable spherical nullhomotopy. This proves
Theorem~\ref{thm:main}.

\section{Exact adversarial local models and certificate scope}
The following low-dimensional models test delicate algebraic cases;
they are \emph{not} numerical representatives of the nonzero class
in \(\pi_9\T\), and are not counterexamples to the manuscript.

First consider
\[
 g_{\mathrm{sing}}(x,y)=-x^2+\frac{y^4}{x^2+y^2}\quad ((x,y)\ne0).
\]
On the unit circle its homotopy from \(q_{1,1}\) has angular
function
\[
 h_t=-x^2+y^2-tx^2y^2,\quad 0\leq t\leq1.
\]
Writing \(s=y^2\), this is \(2s-1-ts(1-s)\); its \(s\)-derivative
is \(2-t+2ts\geq1\), and its endpoint values are \(-1,1\).
There is a unique interior root in \(s\), giving four transverse
angular zeros throughout the homotopy. At \(w=(2,0)\),
\[
 \nabla g_{\mathrm{sing}}=(-4,0),\qquad
 H=\begin{pmatrix}-2&0\\0&0\end{pmatrix}.
\]
For \(R=4,v=(-1,0)\), the vertical Hessian is singular, but the
columns obtained by varying \(w_x\) and the sphere tangent \(v_y\)
are \((-2,0),(0,-4)\), of determinant \(8\).
This verifies the necessity of the full-map regularity criterion.

For a genuine origin-crossing pair of equal vertical gradients put
\[
 g_{\mathrm{cross}}(x,y)=x\sqrt{x^2+y^2}
                       -\frac{2xy^3}{x^2+y^2}.
\]
It is smooth off zero and homogeneous of degree two for positive
scalars. Its angular function is
\(h(\varphi)=\cos\varphi(1-2\sin^3\varphi)\).
Its four zeros are the two zeros of cosine and the two points
with \(\sin\varphi=2^{-1/3}\). At the former its derivatives are
\(1,3\); at the latter they are
\(-6\sin^2\varphi\cos^2\varphi\ne0\).
Thus its negative domain consists of two disjoint arcs with
alternating signs and is in \(T_{1,1}\): an orientation-preserving
circle isotopy matches the ordered zeros and negative intervals
to those of \(q_{1,1}\), and the defining-function fibers are convex.

At \(w=(2,0)\) and \(w'=(-2,0)\), both gradients are \((4,0)\).
With \(R=4,v=(1,0)\), their sphere momenta coincide and their
segment passes through the origin. In a base-direction germ where
\(d_bg_b\) vanishes at those endpoints, their remaining momentum
difference is \(-4\,dR\), nonzero when \(dR\ne0\).
The formula \eqref{eq:mixed}, rather than an assumption that
collision segments avoid zero, handles this case in the general
construction. Such a local angular model can be inserted along a
path from the standard tube and localized in a base disk; its
appearance says nothing about the global stable class.

\medskip
\noindent\textbf{What the certificate certifies.}
\texttt{verify.py} replays \texttt{exact\_certificate.json} using
integer and rational arithmetic and exact symbolic differentiation.
It checks the universal 64-map calculation, including the complete
quotient-type table, and the independent replay of the retained
Jacobian output. It verifies the displayed index/degree
arithmetic, cutoff differentiation, momenta and Euler identities,
the normalized strict-choice margins, and the two explicit local
models. The local model calculations do not replace the
dimension-independent proof of regularity, global inequalities,
compactness, or smooth-family localization above.
The source theorems and their model compatibility are prose and
literature evidence, not Lean/kernel or numerical evidence.
Neither sampling nor a no-collision search is used as proof.
The exact script has no access to, or computed encoding of, the
actual topological class \(a\) or connecting homomorphism.
With Python and SymPy available, the standalone replay from this
directory is \texttt{python3 verify.py --self-test}.

Compared with the strongest prior result in this campaign, the
accepted one-critical obstruction, this establishes its previously
missing geometric input: a nonzero spherically trivial tube class
with the claimed embedded realization. It does not establish that
Hamiltonian transport retains that class or gives the forbidden
one-critical endpoint. Those implications, independent review of
this proof, and campaign-level novelty/publication assessment
remain separate.

\begin{thebibliography}{9}
\bibitem{AGIS}
D. Alvarez-Gavela, K. Igusa and M. G. Sullivan,
\emph{Legendrian and Lagrangian higher torsion},
arXiv:2603.28007v2 (19 August 2026), Section 3.6, Theorem 3.30,
Remark 3.34.
\url{https://arxiv.org/abs/2603.28007v2}.
\bibitem{AAGCK}
M. Abouzaid, D. Alvarez-Gavela, S. Courte and T. Kragh,
\emph{Normal invariant of nearby Lagrangians via twisted derivative},
arXiv:2505.12515v2 (22 October 2025), Definitions 3.10, 3.36,
4.1, 4.3; Propositions 4.4, 4.6--4.8; Remark 4.5.
\url{https://arxiv.org/abs/2505.12515v2}.
\bibitem{Waldhausen}
F. Waldhausen, \emph{Algebraic K-theory of spaces, a manifold approach},
CMS Conference Proceedings 2 (1982), 141--184.
Propositions 3.1--3.2 and Corollary 3.4, retained scan pp.18--20.
The author-hosted original is retained as
\texttt{sources/waldhausen1982.pdf}.
\bibitem{WJR}
F. Waldhausen, B. Jahren and J. Rognes,
\emph{Spaces of PL manifolds and categories of simple maps},
author-hosted prepublication text, Theorems 0.1 and 0.3,
printed pp.2 and 4. Retained as
\texttt{../zero\_band\_audit/wjr2013.pdf} and
\texttt{wjr2013-opening.txt}; the final book appeared in 2013.
\bibitem{Adams}
J. F. Adams, \emph{On the groups \(J(X)\). IV},
Topology 5 (1966), 21--71, Theorem 9.5, p.51.
\url{https://doi.org/10.1016/0040-9383(66)90004-8}.
\bibitem{Rognes}
J. Rognes, \emph{Algebraic K-theory of strict ring spectra},
arXiv:1403.5998v1 (24 March 2014), Theorem 3.3, p.8.
\url{https://arxiv.org/abs/1403.5998v1}.
\bibitem{IWX}
D. C. Isaksen, G. Wang and Z. Xu,
\emph{Stable homotopy groups of spheres: from dimension 0 to 90},
Publications mathematiques de l'IHES 137 (2023), 107--243,
Corollary 1.3, column explanation p.112, Table 1 p.113.
\url{https://doi.org/10.1007/s10240-023-00139-1}.
\end{thebibliography}
\end{document}
